
% Custom tube command
\newcommand{\tube}[1]{%
\begin{tikzpicture}[scale = 0.15]
    \IfSubStr{#1}{a}{ \fill[Blue, draw] (0,4.5) circle (8pt); }{ \draw (0,4.5) circle (8pt); }
    \IfSubStr{#1}{b}{ \fill[Blue, draw] (0,3) circle (8pt); }{ \draw (0,3) circle (8pt); }
    \IfSubStr{#1}{c}{ \fill[Blue, draw] (0,1.5) circle (8pt); }{ \draw (0,1.5) circle (8pt); }
    \IfSubStr{#1}{x}{ \fill[Blue, draw] (1,4.5) circle (8pt); }{ \draw (1,4.5) circle (8pt); }
    \IfSubStr{#1}{y}{ \fill[Blue, draw] (1,3) circle (8pt); }{ \draw (1,3) circle (8pt); }
    \IfSubStr{#1}{z}{ \fill[Blue, draw] (1,1.5) circle (8pt); }{ \draw (1,1.5) circle (8pt); }
\end{tikzpicture}%
}

\begin{tikzpicture}[
    x={(1cm, 0cm)},
    y={(0cm, 0.5cm)},
    z={(0cm, 1cm)},
    scale=.8
]

\def\anglestep{20} % 18 steps * 20 degrees = 360 degree wrap-around

% 1. DRAW BACKGROUND SPHERICAL REFERENCE CIRCLES
\foreach \z/\r in {0/4.47, 2/5.66, 4/6, 6/5.66, 8/4.47}{
    \begin{scope}[canvas is xy plane at z=\z]
        \draw[Orange, thin] (0,0) circle[radius=\r];
    \end{scope}
}

% 2. DEFINE POLES
\coordinate (empty) at (0,0,-2);
\coordinate (abcxyz) at (0,0,10);

% 3. DEFINE EVENLY SPACED COORDINATES 
\begin{scope}[canvas is xy plane at z=0]
    % 6 nodes (spaced by 3)
    \foreach \name/\ax in {a/0, z/3, c/5.5, y/9, b/12, x/15} 
        \coordinate (\name) at (\ax*\anglestep:4.47);
\end{scope}

\begin{scope}[canvas is xy plane at z=2]
    % 9 nodes (spaced by 2)
    \foreach \name/\ax in {ay/0, az/2, bz/4.2, cz/6, cx/9.2, cy/10, by/12, bx/14, ax/15.5} 
        \coordinate (\name) at (\ax*\anglestep:5.66);
\end{scope}

\begin{scope}[canvas is xy plane at z=4]
    % 18 nodes (spaced by 1)
    \foreach \name/\ax in {acy/0, aby/1, ayz/2.2, acz/3, axz/4, abz/5.5, bxz/6, byz/7, bcz/8.1, cyz/9, cxz/10, bcx/11, cxy/12, bcy/13, bxy/14, abx/15, acx/16, axy/17} 
        \coordinate (\name) at (\ax*\anglestep:6);
\end{scope}

\begin{scope}[canvas is xy plane at z=6]
    % 15 nodes (spaced by 1.2)
    \foreach \name/\ax in {abcy/0, abyz/1.2, acyz/2.4, abcz/3.6, acxz/5.2, axyz/6.0, abxz/7.2, bxyz/8.4, bcyz/10.1, cxyz/10.8, bcxz/12.0, bcxy/13.2, abxy/14.4, abcx/15.6, acxy/16.8} 
        \coordinate (\name) at (\ax*\anglestep:5.66);
\end{scope}

\begin{scope}[canvas is xy plane at z=8]
    % 6 nodes (spaced by 3)
    \foreach \name/\ax in {abcyz/2, abcxz/5, acxyz/8, abxyz/10, bcxyz/12.5, abcxy/15} 
        \coordinate (\name) at (\ax*\anglestep:4.47);
\end{scope}


% 4. PRECOMPUTED EDGE LISTS 
\def\backedges{%
    empty/a, empty/c, empty/z,
    z/az, z/bz, z/cz, a/az, c/cz,
    az/ayz, az/acz, az/axz, az/abz,
    bz/abz, bz/bxz, bz/byz, bz/bcz,
    cz/bcz, cz/cyz,
    ay/aby, ay/ayz,
    ayz/abyz, ayz/acyz,
    acz/acyz, acz/abcz, acz/acxz,
    axz/acxz, axz/axyz, axz/abxz,
    abz/abxz, bxz/abxz, bxz/bxyz,
    byz/bxyz, 
    aby/abyz, abyz/abcyz,
    acyz/abcyz, abcz/abcyz, abcz/abcxz,
    axyz/abxyz,byz/bcyz,
    empty/b, empty/x, empty/y,
    acxz/abcxz, acxz/acxyz, axyz/acxyz,c/cx, c/cy,cz/cxz,
    abcy/abcyz%
}

\def\frontedges{%
    abcxyz/acxyz, abcxyz/abxyz, abcxyz/bcxyz, abcxyz/abcxy,abcxyz/abcyz, abcxyz/abcxz,
    a/ax, a/ay, bcz/bcyz,
    y/cy, y/by, b/by, b/bx, x/bx, x/ax,
      aby/abcy,
    ax/abx, ax/acx, ax/axy,
    ay/axy, ay/acy,
    cy/cxy, cy/bcy,
    by/bcy, by/bxy,
    bx/bxy, bx/abx,
    cyz/bcyz, cyz/cxyz,
    cxz/cxyz, cxz/bcxz,
    bcx/bcxz, bcx/bcxy,
    cxy/bcxy, bcy/bcxy, bxy/bcxy, bxy/abxy,
    abx/abxy, abx/abcx,
    acx/abcx, acx/acxy,
    axy/acxy, acy/acxy, acy/abcy,
    abxz/abxyz, bxyz/abxyz, bxyz/bcxyz,
    bcyz/bcxyz, cxyz/bcxyz, bcxz/bcxyz,
    bcxy/bcxyz, bcxy/abcxy,cx/cxz, cx/bcx, cx/cxy,
    abxy/abcxy, abcx/abcxy, acxy/abcxy, abcy/abcxy%
}

% Draw Back Edges (Thin, Gray)
\foreach \u/\v in \backedges {
    \draw[thin, gray] (\u) -- (\v);
}

% Draw Front Edges (Thick, Black)
\foreach \u/\v in \frontedges {
    \draw[thick, black] (\u) -- (\v);
}

% 5. PLACE ALL NODES ON TOP (Removes need for manual redraws later)
\foreach \n in {a,b,c,x,y,z,ax,ay,az,bx,by,bz,cx,cy,cz,abx,aby,abz,acx,acy,acz,bcx,bcy,bcz,axy,bxy,cxy,axz,bxz,cxz,ayz,byz,cyz,abcx,abcy,abcz,abxy,abxz,abyz,acxy,acxz,acyz,bcxy,bcxz,bcyz,axyz,bxyz,cxyz,abcxy,abcxz,abcyz,abxyz,acxyz,bcxyz} {
    \node[fill=white, circle, inner sep=0.5pt] at (\n) {\tube{\n}};
}

% Place Poles
\node[fill=white, circle, inner sep=0.5pt] at (empty) {\tube{}};
\node[fill=white, circle, inner sep=0.5pt] at (abcxyz) {\tube{abcxyz}};

\end{tikzpicture}